\subsection{Quadratic elimination and the repeated lines}
\label{sec:repeatedlines}

The final two leading types have a repeated factor. They reduce
to a polynomial in one variable after eliminating two Gaussian
variables. The full leading-coefficient theorem makes an escape
table unnecessary: only the actual degree, its leading
coefficient and the determinant factor are needed.

\begin{lemma}[The determinant factor in quadratic elimination]
\label{lem:quadraticelimination}
Over $\C(t)$, let a phase have the form
\[
 -\tfrac12 y^T M(t)y+L(t,x)^Ty+V(t,x),
\]
where $M$ is an invertible symmetric matrix and $L,V$ are
polynomial in $x$. Eliminating $y$ identifies its top twisted
module and constant volume seed with the one-variable phase
and normalization
\[
 \varphi=V+\tfrac12L^TM^{-1}L,\qquad
 h=(\det M)^{-1/2}.
\]
The parameter connection is
$\partial_t+\varphi_t+h'/h$. An algebraic square root is allowed.
If the actual degree $d$ of $\varphi$ is at least two and
$d\,v_\infty(h)>v_\infty([x^d]\varphi)$, its constant seed
cannot have a rational unit quotient.
\end{lemma}
\begin{proof}
The change $y=z+M^{-1}L$ has Jacobian one. Its spatial twisted
relations reduce the $z$ factor to the one-dimensional Gaussian
quotient, with normalized contractions
$\ell(z_i z_j)=(M^{-1})_{ij}$. Parameter differentiation of its
constant class inserts $-z^TM'z/2$, hence
\[
 -\tfrac12\operatorname{tr}(M^{-1}M')=h'/h.
\]
Terms from the moving center are spatial exact forms, by the
product rule or Cartan's formula. This proves the identification
of both connection and seed. After a finite cover at infinity,
$h$ is a Laurent series and Theorem~\ref{thm:stationaryballs}
applies. The cover multiplies both sides of the valuation
inequality by the same positive integer.

The elimination is over the generic field. It need not define
a conditional formal average at $t=0$. The original Wick
functional first gives its rational seed quotient by
Lemma~\ref{lem:seedfunctional}; this quotient is transported
through the displayed module identification. No image-kernel
calculation or convergent original period is needed.
\end{proof}

\begin{lemma}\label{lem:quadraticfibrehighdegree}
Let $u,Y_1,\ldots,Y_s$ be independent standard Gaussian variables
and
\[
 P=Y^TAY+B(u)^TY+g(u),\qquad A=A^T,
\]
with $B,g$ arbitrary polynomials. Put
\[
 D(t)=\det(I-2tA),\quad C(t)=(I-2tA)^{-1},\quad
 \varphi=tg-u^2/2+(t^2/2)B^TCB.
\]
If $\deg_u\varphi\ge3$, then $P$ is not pure.
\end{lemma}
\begin{proof}
Quadratic elimination gives $h=D^{-1/2}$. Let $d=\deg D$ and
$r=\deg_u\varphi\ge3$. Its leading coefficient is $N(t)/D(t)$,
where $N$ is a nonzero polynomial divisible by $t$:
the degree-$r$ coefficient vanishes at $t=0$. Consequently
\[
 r v_\infty(h)-v_\infty(N/D)
     =(r/2-1)d+\deg N>0.
\]
Apply Lemma~\ref{lem:quadraticelimination}. This estimate is
independent of the degrees of $B,g$ and the dimension $s$.
\end{proof}

\subsubsection*{A nonisotropic repeated factor}

Normalize the repeated factor to $u$ with $\E[u^2]=1$ and take
an orthogonal Gaussian complement $Y\in\C^2$. Every such cubic
has the form in Lemma~\ref{lem:quadraticfibrehighdegree}, with
\[
 B(u)=p u^2+q u+r,\qquad\deg g\le3.
\]
For a double line $p\ne0$. If its effective quartic coefficient
is nonzero, the lemma excludes purity immediately. If that
coefficient vanishes identically, then
\[
 0=p^T\operatorname{adj}(I-2tA)p
   =p^Tp-2t\bigl(\operatorname{tr}(A)p^Tp-p^TAp\bigr).
\]
Thus $p^Tp=p^TAp=0$. The nonzero vector $(0,p)$ is isotropic,
lies in the leading singular plane $u=0$, and satisfies
$F_2(0,p)=0$. It supplies the frame of Lemma~\ref{lem:frame}.
For a triple line, $p=0$ and $[u^3]g\ne0$. The effective
phase has exact cubic coefficient $t[u^3]g$, so the same
high-degree lemma excludes purity. This exhausts all matrices
$A$, including nilpotent and defective ones, without a matrix
case split.

\subsubsection*{An isotropic repeated factor}

Use Gaussian circular coordinates with repeated factor $Z$.
The singular plane is $Z=0$, whose only isotropic direction is
the $W$ axis. Write $\lambda=[W^2]F_2$. If $\lambda=0$, the
frame test applies. Assume $\lambda\ne0$.

For a double line write its other factor as
$aW+bT+cZ$. The Gaussian null rotation
\[
 Z=Z',\quad T=T'+sZ',\quad W=W'-sT'-s^2Z'/2
\]
removes the $T$ term if $a\ne0$, or removes the $Z$ term if
$a=0,b\ne0$. Scaling $P$ gives exactly these leading forms:
\[
 \mathrm I: Z^2W+a_0Z^3,\qquad
 \mathrm {II}:Z^2T,\qquad \mathrm {III}:Z^3.
\]
The last is the triple line. The rotation preserves the Gaussian
and keeps the $W^2$ coefficient nonzero. All lower terms remain
arbitrary.

Write $P=\lambda W^2+W H(Z,T)+G(Z,T)$. In type I,
$H=Z^2+bZ+cT+d$ and $G=a_0Z^3+G_2$. In types II and III,
$H=bZ+cT+d$ and $G=Z^2T+G_2$ or $Z^3+G_2$ respectively.
Here $G_2$ is an arbitrary polynomial of degree at most two.
Eliminating $W$ in $tP-ZW-T^2/2$ gives
\[
 \Psi=t\left(G-\frac{H^2}{4\lambda}\right)
       +\frac{ZH}{2\lambda}-\frac{T^2}{2}
       -\frac{Z^2}{4\lambda t}.
\]
Its $T^2$ coefficient is $C_T=a t-1/2$, where
$a=g_{02}-c^2/(4\lambda)$ and $g_{02}=[T^2]G_2$.
Eliminate $T$ as well. If $e=\deg_t C_T\in\{0,1\}$, the total
normalization has $v_\infty(h)=(1+e)/2$. The actual remaining
leading coefficients are
\[
 \begin{array}{c|c|c}
 \text{type}&[Z^d]\varphi&d\\\hline
 \mathrm I&\displaystyle
  -\frac{t}{4\lambda}-\frac{c^2t^2}{16\lambda^2C_T}
   =-\frac{t(2g_{02}t-1)}{8\lambda C_T}&4\\[5pt]
 \mathrm {II}&\displaystyle-\frac{t^2}{4C_T}&4\\[4pt]
 \mathrm {III}&t&3.
 \end{array}
\]
None is identically zero. Their valuations are at most $e-1$,
$e-2$, and $-1$, respectively. Each is strictly less than
$d(1+e)/2$. Lemma~\ref{lem:quadraticelimination} excludes
purity in all three types. In particular this includes the
boundary where the quartic coefficient tends to a nonzero
constant while lower coefficients grow.

\begin{proposition}\label{prop:repeatedcomplete}
Every pure ternary cubic with repeated leading factor has a
Gaussian affine-linear coordinate. It therefore satisfies
$m>3\deg Q$.
\end{proposition}
\begin{proof}
The nonisotropic and isotropic alternatives above exhaust the
repeated factor. In every surviving branch the explicit frame
test succeeds. Apply the complete affine-linear theorem.
\end{proof}
